\section{Reti di telefonia mobile}
All'inizio, quindi anni 90, abbiamo ogni 10 anni una nuova generazione di reti mobili, con un salto tecnologico importante. Oggi siamo alla 5G, ma è già in fase di sviluppo la 6G.

TACS (Total Access Communication System) è stata la prima generazione di reti mobili, introdotta negli anni 80. Era una rete analogica, che permetteva solo chiamate vocali e aveva una capacità limitata.
AMPS (Advanced Mobile Phone System) è stata la prima rete mobile digitale, introdotta negli anni 80. Permetteva chiamate vocali e messaggi di testo, ma aveva ancora una capacità limitata.\\

Negli anni 90 sono stati introdotti i primi standard di rete mobile digitale, come GSM (Global System for Mobile Communications) con cui abbiamo avuto connessioni digitali, avevamo sia timeslot che frequenza.
EDGE (Enhanced Data rates for GSM Evolution) è stato un miglioramento del GSM, introdotto negli anni 2000. Permetteva velocità di trasmissione dati più elevate, fino a 384 kbps, ma era ancora limitato rispetto alle reti successive.
GPRS (General Packet Radio Service) è stato un miglioramento del GSM. Permetteva la trasmissione di dati a pacchetto, con velocità fino a 114 kbps, ma era ancora limitato rispetto alle reti successive.

Quando abbiamo introdotto i sistemi GPRS a EDGE siamo passati dalla seconda generazione alla 2.5.\\
La terza generazione fa riferiferimento a UMTS (Universal Mobile Telecommunications System), introdotto negli anni 2000. Permetteva velocità di trasmissione dati più elevate, fino a 2 Mbps, e supportava servizi multimediali come videochiamate e streaming.
Si assegnano codici e non piú timeslot o frequenze.
HSPA (High Speed Packet Access) è stato un miglioramento del UMTS, introdotto negli anni 2000. Permetteva velocità di trasmissione dati più elevate, fino a 14 Mbps, ma era ancora limitato rispetto alle reti successive.
Si passa da una tecnica ibrida time division multiple access (TDMA) e code division multiple access (CDMA) a una tecnica completamente basata su CDMA.\\
La quarta generazione fa riferimento a LTE (Long Term Evolution), introdotto negli anni 2010. Permetteva velocità di trasmissione dati più elevate, fino a 100 Mbps, e supportava servizi multimediali avanzati come video in alta definizione e giochi online.
LTE non é i sistemi 4G, per dire che é un sistema 4G allora stiamo parlando di LTE-Advanced, che permette velocità di trasmissione dati fino a 1 Gbps.\\
La quinta generazione fa riferimento a 5G, introdotto negli anni 2020. Permette velocità di trasmissione dati molto elevate, fino a 10 Gbps, e supporta servizi multimediali avanzati come realtà virtuale e aumentata, e l'Internet of Things (IoT).
Il 5G si basa su 3 pilastri, l'EMBB (Enhanced Mobile Broadband) per la connettività ad alta velocità, l'URLLC (Ultra-Reliable Low Latency Communications) per applicazioni che richiedono bassa latenza e alta affidabilità, e mMTC (massive Machine Type Communications) per supportare un gran numero di dispositivi connessi, come quelli dell'IoT.
Il 5G non ha un vero nome ma lo si chiama New Radio (NR).\\

I sistemi 6G non sono ancora stati standardizzati, ma si prevede che offriranno velocità di trasmissione dati ancora più elevate, fino a 100 Gbps, e supportino servizi multimediali avanzati come realtà virtuale e aumentata, e l'Internet of Things (IoT) in modo ancora più efficiente. Si prevede che i sistemi 6G utilizzeranno tecnologie come l'intelligenza artificiale e la comunicazione quantistica per migliorare ulteriormente le prestazioni delle reti mobili.

\subsection{1G}
Martin Cooper inventore del primo telefono cellulare, ha effettuato la prima chiamata da un telefono cellulare il 3 aprile 1973. Il telefono utilizzato era un Motorola DynaTAC, che pesava circa 1 kg e aveva un'autonomia di circa 30 minuti in conversazione.
In realtá i primi cellulari erano disponibili per le auto che si collegavano alla batteria, infatti il problema principale era creare un telefono con autonomia di batteria.
Si chiama cellulare perché la copertura della rete è suddivisa in celle, ognuna delle quali è servita da una stazione base. Quando un utente si sposta da una cella all'altra, la chiamata viene trasferita automaticamente alla nuova stazione base, garantendo una continuità di servizio.
Le celle avevano una dimensione fissa di circa 10 km, e la capacità di ogni cella era limitata a circa 100 utenti. Questo ha portato a problemi di congestione nelle aree urbane, dove la domanda di servizi mobili era elevata.\\
La posizione dell'utente all'interno delle celle era gestita da un sistema piú o meno automatico, che utilizzava segnali radio per determinare la posizione dell'utente e instradare le chiamate in modo efficiente.
L'assegnazione delle varie frequenze era gestita da un sistema di controllo centrale, che assegnava le frequenze alle stazioni base in base alla domanda e alla disponibilità. Questo sistema era basato su un algoritmo di assegnazione delle frequenze che cercava di minimizzare l'interferenza tra le celle adiacenti. Quando io mi sposto da una cella all'altra, il sistema di controllo centrale gestisce il processo di handover, che consiste nel trasferire la chiamata dalla stazione base della cella precedente a quella della nuova cella. Questo processo è gestito in modo da garantire una continuità di servizio e minimizzare l'interruzione della chiamata.\\
Quando questo sistemo é stato sviluppato negli anni 80, l'apparecchio e fare o ricevere le telefonate era molto costoso, per questo si limitavano le chiamate a quelle più importanti, e si preferiva utilizzare i telefoni fissi per le chiamate quotidiane. Inoltre, la copertura della rete era limitata, soprattutto nelle aree rurali, il che ha ulteriormente limitato l'adozione dei telefoni cellulari.
Il sistema é morto perché lo hanno clonato, quando l'utente trasmetteva e si agganciava trasmetteva in chiaro il suo numero di telefono e sia il numero identificativo del telefono e questo causava problemi di sicurezza, infatti era possibile clonare il telefono e fare chiamate a spese dell'utente originale. 
\subsection{2G}
La seconda generazione di reti mobili, introdotta negli anni 90, ha portato a una rivoluzione nella comunicazione mobile. Con l'introduzione del GSM (Global System for Mobile Communications), le chiamate vocali sono diventate digitali, migliorando la qualità del suono e la sicurezza delle comunicazioni. Inoltre, il GSM ha introdotto la possibilità di inviare messaggi di testo (SMS), che è diventato un modo popolare per comunicare brevemente e rapidamente. La capacità della rete è aumentata significativamente, consentendo a più utenti di connettersi contemporaneamente senza problemi di congestione. Tuttavia, il GSM aveva ancora limitazioni in termini di velocità di trasmissione dati, che sono state superate con l'introduzione di tecnologie come GPRS ed EDGE, che hanno portato alla transizione verso la 2.5G.\\
Hanno inventato la possibilitá di identificare il mobile terminal con la SIM (Subscriber Identity Module), che è una scheda rimovibile che contiene informazioni sull'abbonato, come il numero di telefono e le credenziali di autenticazione. Questo ha permesso agli utenti di cambiare telefono mantenendo lo stesso numero di telefono e ha migliorato la sicurezza delle comunicazioni, poiché le informazioni dell'abbonato erano memorizzate in modo sicuro sulla SIM.
La tecnica di modulazione era Gaussian Minimum Shift Keying (GMSK), che è una tecnica di modulazione digitale che utilizza una forma d'onda gaussiana per modulare il segnale. Questa tecnica è stata scelta per il GSM perché offre una buona efficienza spettrale e una bassa interferenza tra i canali adiacenti, consentendo di utilizzare le frequenze in modo più efficiente e migliorando la qualità delle comunicazioni.\\
La trasmissione era digitale e usava insieme sia frequency division multiple access (FDMA) che time division multiple access (TDMA). FDMA divide lo spettro di frequenza in canali separati, mentre TDMA divide il tempo in slot, consentendo a più utenti di condividere la stessa frequenza. Se uso TDMA significa che devo avere un aggancio sincronizzato, se invece uso FDMA non ho bisogno di un aggancio sincronizzato, ma ho bisogno di un filtro più stretto per evitare interferenze tra i canali adiacenti.
La tolleranza che potevo accettare era di 100 microsecondi, che è il tempo massimo che potevo tollerare per la sincronizzazione tra i terminali e la stazione base. Se la tolleranza fosse stata maggiore, avrei avuto problemi di interferenza tra i canali adiacenti, mentre se fosse stata minore, avrei avuto problemi di sincronizzazione e di perdita di pacchetti.\\

\paragraph{TDMA-FDMA}
Si usavano dei canali e il GSM era chiamato dual band perché potevo usare sia la banda a 900 MHz che quella a 1800 MHz. Ogni canale aveva una larghezza di banda di 200 kHz, e ogni canale era diviso in 8 timeslot, ognuno dei quali poteva essere assegnato a un utente diverso. In questo modo, ogni canale poteva supportare fino a 8 utenti contemporaneamente, e la capacità totale della rete era aumentata significativamente rispetto alla prima generazione.\\
\paragraph{Architettura di una rete 2G}
L'architettura di una rete 2G é rimasta. Il cellulare viene chiamato mobile terminal o mobile equipment, al suo interno ha la SIM che identifica l'utente. La connesione tra la Base station e la stazione base controller è chiamata Abis, mentre la connessione tra la stazione base controller e il mobile switching center è chiamata A. Il mobile switching center è responsabile della gestione delle chiamate e della connessione alla rete fissa, mentre la stazione base controller gestisce le risorse radio e la comunicazione con le stazioni base.\\
La Base station viene controllata da un BSC (Base Station Controller) che gestisce le risorse radio e la comunicazione con le stazioni base, mentre il MSC (Mobile Switching Center) è responsabile della gestione delle chiamate e della connessione alla rete fissa.\\
Devo avere informazioni relative all'home register e al visitor register, che sono database che contengono informazioni sugli utenti della rete. L'home register contiene informazioni sugli utenti registrati nella rete, mentre il visitor register contiene informazioni sugli utenti che si trovano temporaneamente nella rete, ad esempio quando si spostano da una cella all'altra. Queste informazioni sono utilizzate per gestire le chiamate e garantire la continuità del servizio quando gli utenti si spostano tra le celle.\\
EIR é l'informazione non sull'utente ma sul dispositivo che ha un database che contiene informazioni sui dispositivi mobili, come il numero di serie e il modello. Questo database viene utilizzato per identificare i dispositivi mobili e garantire la sicurezza della rete, ad esempio bloccando i dispositivi rubati o non autorizzati.\\

La SIm ha codice di sicurezza personali come PIN (Personal Identification Number) e PUK (Personal Unblocking Key) che sono utilizzati per proteggere l'accesso alla SIM e garantire la sicurezza delle comunicazioni. Il PIN è un codice a 4 cifre che l'utente deve inserire per accedere alla SIM, mentre il PUK è un codice a 8 cifre che viene utilizzato per sbloccare la SIM se il PIN viene inserito in modo errato più volte. Questi codici sono importanti per proteggere le informazioni dell'utente e garantire la sicurezza della rete.\\
IMEI é un codice univoco che identifica ogni dispositivo mobile, ed è utilizzato per garantire la sicurezza della rete e prevenire l'uso di dispositivi non autorizzati. Il codice IMEI viene assegnato a ogni dispositivo mobile durante la produzione e viene memorizzato nella SIM. Quando un dispositivo si connette alla rete, il codice IMEI viene verificato per garantire che il dispositivo sia autorizzato a utilizzare la rete. Se un dispositivo viene segnalato come rubato o non autorizzato, il codice IMEI può essere bloccato, impedendo al dispositivo di connettersi alla rete.\\

\paragraph{General Packet Radio Service - GPRS}
GPRS é un miglioramento del GSM che permette la trasmissione di dati a pacchetto, con velocità fino a 114 kbps. GPRS utilizza una tecnica di modulazione chiamata Gaussian Minimum Shift Keying (GMSK) e supporta sia la trasmissione dati che la trasmissione vocale. GPRS è stato un passo importante verso la 3G, poiché ha introdotto la possibilità di trasmettere dati a pacchetto, consentendo l'accesso a Internet e ad altri servizi basati sui dati. Tuttavia, GPRS aveva ancora limitazioni in termini di velocità di trasmissione dati, che sono state superate con l'introduzione di tecnologie come EDGE e UMTS.\\
Arriva da un internet service provider un indirizzo IP e il gateway deve inoltrare il pacchetto verso il base station. 

\subsection{3G}
La terza generazione di reti mobili, introdotta negli anni 2000, ha portato a un ulteriore miglioramento nella comunicazione mobile. Con l'introduzione del UMTS (Universal Mobile Telecommunications System), le velocità di trasmissione dati sono aumentate significativamente, consentendo l'accesso a Internet e ad altri servizi basati sui dati in modo più efficiente. Inoltre, il UMTS ha introdotto la possibilità di supportare servizi multimediali come videochiamate e streaming, migliorando ulteriormente l'esperienza dell'utente. Tuttavia, il UMTS aveva ancora limitazioni in termini di velocità di trasmissione dati, che sono state superate con l'introduzione di tecnologie come HSPA e LTE, che hanno portato alla transizione verso la 4G.\\
In italia il primo operatore a lanciare la rete 3G è stato 3 Italia, seguito da Vodafone e TIM. La copertura della rete 3G è stata inizialmente limitata alle aree urbane, ma si è rapidamente estesa a livello nazionale, consentendo a un numero sempre maggiore di utenti di accedere ai servizi basati sui dati. La rete 3G ha rappresentato un passo importante verso la connettività mobile avanzata, aprendo la strada a nuove applicazioni e servizi che hanno trasformato il modo in cui le persone comunicano e accedono alle informazioni.\\
\paragraph{Wideband Code Division Multiple Access - WCDMA}
WCDMA é una tecnica di accesso multiplo utilizzata nelle reti 3G, che consente a più utenti di condividere la stessa frequenza utilizzando codici univoci per distinguere i segnali. WCDMA utilizza una tecnica di modulazione chiamata Quadrature Phase Shift Keying (QPSK) e supporta sia la trasmissione dati che la trasmissione vocale. WCDMA è stato un passo importante verso la 4G, poiché ha introdotto la possibilità di trasmettere dati a pacchetto in modo più efficiente, consentendo l'accesso a Internet e ad altri servizi basati sui dati in modo più veloce e affidabile. Tuttavia, WCDMA aveva ancora limitazioni in termini di velocità di trasmissione dati, che sono state superate con l'introduzione di tecnologie come HSPA e LTE.\\
\paragraph{High Speed Packet Access - HSPA}
HSPA é un miglioramento del UMTS che permette velocità di trasmissione dati più elevate, fino a 14 Mbps. HSPA utilizza una tecnica di modulazione chiamata Quadrature Phase Shift Keying (QPSK) e supporta sia la trasmissione dati che la trasmissione vocale. HSPA è stato un passo importante verso la 4G, poiché ha introdotto la possibilità di trasmettere dati a pacchetto in modo più efficiente, consentendo l'accesso a Internet e ad altri servizi basati sui dati in modo più veloce e affidabile. Tuttavia, HSPA aveva ancora limitazioni in termini di velocità di trasmissione dati, che sono state superate con l'introduzione di tecnologie come LTE, che hanno portato alla transizione verso la 4G.\\
\paragraph{Architettura di una rete 3G}
L'architettura di una rete 3G è simile a quella di una rete 2G, ma con alcune modifiche per supportare le nuove tecnologie e i nuovi servizi. La rete 3G utilizza ancora stazioni base e stazioni base controller, ma introduce anche nuovi componenti come il Node B, che è la stazione base utilizzata nelle reti 3G, e il Radio Network Controller (RNC), che gestisce le risorse radio e la comunicazione tra i Node B. Inoltre, la rete 3G introduce un nuovo componente chiamato Serving GPRS Support Node (SGSN), che gestisce la trasmissione dei dati a pacchetto e la connessione alla rete fissa. L'architettura di una rete 3G è progettata per supportare servizi multimediali avanzati e velocità di trasmissione dati più elevate rispetto alle reti precedenti.\\
Lo User Equipment é il tewrminale dell'utente. Ora c'é la USIM (Universal Subscriber Identity Module) che è una versione avanzata della SIM, che supporta funzionalità aggiuntive come la crittografia e l'autenticazione avanzata. La USIM è utilizzata nelle reti 3G e successive per garantire la sicurezza delle comunicazioni e proteggere le informazioni dell'utente.\\
La RAN (Radio Access Network) è la parte della rete che gestisce la comunicazione tra i terminali mobili e le stazioni base. La RAN è composta da componenti come il Node B e il Radio Network Controller (RNC), che gestiscono le risorse radio e la comunicazione tra i terminali mobili e le stazioni base. La RAN è responsabile della gestione delle risorse radio, della trasmissione dei dati a pacchetto e della connessione alla rete fissa, garantendo una comunicazione efficiente e affidabile per gli utenti mobili.\\
IL nodo B converte il flusso dei dati digitali in segnali radio usando la modulazione WCDMA. \\
La parte di rete Core si divide in due parti, la parte di rete di commutazione (Circuit Switched - CS) che gestisce le chiamate vocali e la parte di rete a pacchetto (Packet Switched - PS) che gestisce la trasmissione dei dati a pacchetto. La parte di rete CS è responsabile della gestione delle chiamate vocali, mentre la parte di rete PS è responsabile della gestione della trasmissione dei dati a pacchetto e della connessione alla rete fissa.\\
I registri comuni sono rimasti piú o meno gli stessi, abbiamo ancora l'home register e il visitor register, ma abbiamo anche un nuovo registro chiamato Authentication Center (AuC), che è responsabile della gestione delle credenziali di autenticazione degli utenti e della generazione dei codici di autenticazione per garantire la sicurezza delle comunicazioni.\\

\paragraph{Multiple Input Multiple Output}
MIMO é una tecnologia utilizzata nelle reti 3G e successive che consente di utilizzare più antenne per trasmettere e ricevere segnali, migliorando la qualità del segnale e aumentando la capacità della rete. MIMO utilizza tecniche di codifica spaziale per trasmettere più flussi di dati contemporaneamente, consentendo di aumentare la velocità di trasmissione dati e migliorare l'affidabilità delle comunicazioni. MIMO è stato un passo importante verso la 4G, poiché ha introdotto la possibilità di trasmettere dati a pacchetto in modo più efficiente, consentendo l'accesso a Internet e ad altri servizi basati sui dati in modo più veloce e affidabile. Tuttavia, MIMO aveva ancora limitazioni in termini di velocità di trasmissione dati, che sono state superate con l'introduzione di tecnologie come LTE, che hanno portato alla transizione verso la 4G.\\
Tre tipi di MIMO: Spatioal Multiplexing, Diversity Coding e Beamforming. Il Spatial Multiplexing consente di trasmettere più flussi di dati contemporaneamente utilizzando più antenne, aumentando la velocità di trasmissione dati. Il Diversity Coding utilizza più antenne per trasmettere lo stesso flusso di dati, migliorando l'affidabilità delle comunicazioni in condizioni di segnale debole o interferenze. Il Beamforming utilizza più antenne per focalizzare il segnale verso un utente specifico, migliorando la qualità del segnale e aumentando la capacità della rete.\\

\subsection{Long Term Evolution}
La quarta generazione di reti mobili, introdotta negli anni 2010, ha portato a un ulteriore miglioramento nella comunicazione mobile. Con l'introduzione del LTE (Long Term Evolution), le velocità di trasmissione dati sono aumentate significativamente, consentendo l'accesso a Internet e ad altri servizi basati sui dati in modo più efficiente. Inoltre, il LTE ha introdotto la possibilità di supportare servizi multimediali avanzati come video in alta definizione e giochi online, migliorando ulteriormente l'esperienza dell'utente. Tuttavia, il LTE aveva ancora limitazioni in termini di velocità di trasmissione dati, che sono state superate con l'introduzione di tecnologie come 5G, che hanno portato alla transizione verso la 5G.\\
Le caratteristiche principali sono la mobilitá, sicurezza, roaming e prestazioni.
Per la mobilitá abbiamo la possibilità di spostarsi da una cella all'altra senza interruzioni, grazie al processo di handover che consente di trasferire la chiamata o la connessione dati da una stazione base a un'altra in modo fluido e senza interruzioni.\\
Per la sicurezza, il LTE utilizza meccanismi di autenticazione e crittografia avanzati per proteggere le comunicazioni e garantire la privacy degli utenti.\\
Per il roaming, il LTE supporta la possibilità di utilizzare i servizi mobili anche quando si è
fuori dalla propria rete domestica, consentendo agli utenti di connettersi a reti partner in tutto il mondo.\\
Per le prestazioni, il LTE offre velocità di trasmissione dati elevate, bassa latenza e una maggiore capacità della rete, consentendo agli utenti di accedere a servizi basati sui dati in modo più efficiente e affidabile.\\

MIMO é una tecnologia utilizzata nelle reti 4G e successive che consente di utilizzare più antenne per trasmettere e ricevere segnali, migliorando la qualità del segnale e aumentando la capacità della rete. MIMO utilizza tecniche di codifica spaziale per trasmettere più flussi di dati contemporaneamente, consentendo di aumentare la velocità di trasmissione dati e migliorare l'affidabilità delle comunicazioni. MIMO è stato un passo importante verso la 5G, poiché ha introdotto la possibilità di trasmettere dati a pacchetto in modo più efficiente, consentendo l'accesso a Internet e ad altri servizi basati sui dati in modo più veloce e affidabile. Tuttavia, MIMO aveva ancora limitazioni in termini di velocità di trasmissione dati, che sono state superate con l'introduzione di tecnologie come 5G, che hanno portato alla transizione verso la 5G.\\
SC-OFDMA é una tecnica di accesso multiplo utilizzata nelle reti 4G, che consente a più utenti di condividere la stessa frequenza utilizzando codici univoci per distinguere i segnali. SC-OFDM utilizza una tecnica di modulazione chiamata Orthogonal Frequency Division Multiplexing (OFDM) e supporta sia la trasmissione dati che la trasmissione vocale. SC-OFDM è stato un passo importante verso la 5G, poiché ha introdotto la possibilità di trasmettere dati a pacchetto in modo più efficiente, consentendo l'accesso a Internet e ad altri servizi basati sui dati in modo più veloce e affidabile. Tuttavia, SC-OFDM aveva ancora limitazioni in termini di velocità di trasmissione dati, che sono state superate con l'introduzione di tecnologie come 5G, che hanno portato alla transizione verso la 5G.\\
SDR (Software Defined Radio) é una tecnologia utilizzata nelle reti 4G e successive che consente di implementare le funzionalità radio in software, invece che in hardware. SDR consente di aggiornare e modificare le funzionalità radio in modo più flessibile e rapido, consentendo di adattarsi alle nuove tecnologie e ai nuovi standard senza dover sostituire l'hardware. SDR è stato un passo importante verso la 5G, poiché ha introdotto la possibilità di implementare nuove funzionalità radio in modo più efficiente, consentendo l'accesso a Internet e ad altri servizi basati sui dati in modo più veloce e affidabile. Tuttavia, SDR aveva ancora limitazioni in termini di velocità di trasmissione dati, che sono state superate con l'introduzione di tecnologie come 5G, che hanno portato alla transizione verso la 5G.\\
Scheduling Dinamico é una tecnica utilizzata nelle reti 4G e successive che consente di allocare le risorse radio in modo dinamico in base alle esigenze degli utenti e alle condizioni della rete. Lo scheduling dinamico utilizza algoritmi di scheduling per determinare quali utenti devono essere serviti in un dato momento, tenendo conto di fattori come la qualità del segnale, la quantità di dati da trasmettere e la priorità degli utenti. Lo scheduling dinamico è stato un passo importante verso la 5G, poiché ha introdotto la possibilità di allocare le risorse radio in modo più efficiente, consentendo l'accesso a Internet e ad altri servizi basati sui dati in modo più veloce e affidabile. Tuttavia, lo scheduling dinamico aveva ancora limitazioni in termini di velocità di trasmissione dati, che sono state superate con l'introduzione di tecnologie come 5G, che hanno portato alla transizione verso la 5G.\\
\paragraph{OFDM e OFDMA}
OFDM (Orthogonal Frequency Division Multiplexing) é una tecnica di modulazione utilizzata nelle reti 4G e successive che consente di trasmettere dati su più frequenze contemporaneamente, migliorando la velocità di trasmissione dati e l'efficienza spettrale. OFDM divide il canale di comunicazione in sottocanali ortogonali, consentendo di trasmettere più flussi di dati contemporaneamente senza interferenze. OFDMA (Orthogonal Frequency Division Multiple Access) é una tecnica di accesso multiplo basata su OFDM che consente a più utenti di condividere la stessa frequenza utilizzando codici univoci per distinguere i segnali. OFDMA è stato un passo importante verso la 5G, poiché ha introdotto la possibilità di trasmettere dati a pacchetto in modo più efficiente, consentendo l'accesso a Internet e ad altri servizi basati sui dati in modo più veloce e affidabile. Tuttavia, OFDMA aveva ancora limitazioni in termini di velocità di trasmissione dati, che sono state superate con l'introduzione di tecnologie come 5G, che hanno portato alla transizione verso la 5G.\\

